\documentclass[11pt]{article}

\usepackage[margin=1in]{geometry}
\usepackage{amsmath,amssymb,amsthm}
\usepackage{array,booktabs,float,tabularx}
\usepackage[round,authoryear]{natbib}
\usepackage{hyperref}

\title{Internalizing Common-Noun Identity Criteria with Quotient Higher Inductive Types}
\author{Anonymous manuscript}
\date{2026-07-10}

\newcommand{\CN}{\mathsf{CN}}
\newcommand{\Raw}{\mathsf{Raw}}
\newcommand{\IC}{\mathsf{IC}}
\newcommand{\Quot}[2]{#1/{#2}}
\newcommand{\embed}[1]{\lvert #1\rvert}
\newcommand{\sem}[1]{[\![#1]\!]}
\newcommand{\Prop}{\mathsf{Prop}}
\newcommand{\keywords}[1]{\par\medskip\noindent\textbf{Keywords:} #1}
\newtheorem{theorem}{Theorem}

\begin{document}
\maketitle

\begin{abstract}
Common-noun semantics needs criteria of identity as well as criteria of
application.  We examine whether a common noun first represented as a setoid
$(\Raw_N,\IC_N)$ can instead be presented as the quotient higher inductive type
$\Quot{\Raw_N}{\IC_N}$, with its identity criterion internalized as equality
paths.  The contribution is deliberately bounded.  We specify the
non-dependent quotient interface required for identity-respecting,
set-valued observations; prove a conditional converse representation theorem;
and provide a Rocq-checked fragment whose global assumptions are explicitly
audited.  The theorem requires function extensionality, proof irrelevance for
respectfulness witnesses, and proposition-valued quotient induction.  Three
linguistic tests connect the interface to semantic concerns: passenger identity
is finer than person identity, modifier evidence need not create extra noun
instances, and physical and informational observations of a book descend
separately while book identity remains conjunctive.  The paper does not claim a
native quotient-HIT implementation in Rocq or a general theory of dot-type
coercions.  It offers a reproducible semantic application of quotient-HIT
ideas with precise formal boundaries.
\end{abstract}

\keywords{common nouns as types, identity criteria, quotient higher inductive types, type-theoretical semantics, Rocq formalization, copredication}

\section{Introduction}

The CN-as-types programme is motivated by the idea that common nouns carry not
only criteria of application but also criteria of identity.  A predicate may
classify an object as a passenger or a book; a common noun analysis must also
say when two such objects count as the same passenger or the same book.  Setoid
semantics expresses this point directly by assigning a carrier together with
an equivalence relation.  That representation is precise, but it creates a
bookkeeping burden: every semantic operation must be paired with a proof that
it respects the relevant equivalence relation.

The proposed transition, CN-as-setoids to CNs-as-quotient-HITs, replaces
this external bookkeeping by a quotient-HIT presentation.  For a common noun
$N$, let $\Raw_N$ be the raw carrier and let
$\IC_N:\Raw_N\to\Raw_N\to\Prop$ be the identity criterion.  The intended
semantic object is:
\[
  \sem{N} = \Quot{\Raw_N}{\IC_N}.
\]
The contribution is not a claim that all of homotopy type theory should be
imported into lexical semantics.  It is a focused claim about quotient HITs as
a way to make identity criteria internal to common-noun types.

\section{Related Work and Position}

The linguistic starting point is the CN-as-types programme, in which a common
noun has both a criterion of application and a criterion of identity
\cite{Luo2012CN}.  Subsequent work on individuation criteria and dot-types
makes the pressure on the identity component especially clear for modification
and copredication \cite{ChatzikyriakidisLuo2015,ChatzikyriakidisLuo2018}.
Our proposal does not replace that analysis.  It asks a narrower question:
given a setoid presentation already motivated by the semantics, can its
identity relation be represented as equality paths in a quotient common-noun
type?

On the type-theoretic side, quotient types are a standard higher-inductive
construction \cite{HoTTBook}, and general semantic accounts of higher
inductive types explain why their constructors and eliminators may be treated
as a coherent interface \cite{LumsdaineShulman2017}.  Cubical work supplies an
important computational setting for HITs \cite{CoquandHuberMortberg2018}.
This paper does not contribute a new HIT implementation or a new metatheory of
quotients.  Its contribution is instead a scoped semantic application: it
identifies the quotient interface needed by common-noun identity criteria,
separates the conditional converse from a stronger unproved claim, and checks
the resulting fragment in Rocq.

More specifically, the contribution is neither a replacement for the setoid
analysis nor a re-description of quotient types in linguistic vocabulary.
The setoid analysis supplies the noun-specific identity criterion; the present
proposal specifies how that criterion is internalized as paths and how
identity-respecting observations descend.  The dot-type literature supplies
the motivation for distinct physical and informational observations of a book;
the present fragment isolates the quotient-recursion obligations that license
those observations.  Finally, the type-theoretic literature supplies quotient
constructors and eliminators; this paper states the exact non-dependent,
set-valued interface used by the semantics and proves a conditional converse
with its assumptions exposed.  These three links are tested by the
passenger/person-times, modifier-evidence, and book/copredication cases below,
rather than asserted as a general theory of coercion or dependent elimination.

\section{Formal Interface}

The quotient presentation has three constructors or structural principles:
\[
\begin{array}{rcl}
  \embed{-} &:& \Raw_N \to \Quot{\Raw_N}{\IC_N},\\
  \mathsf{glue} &:&
    \prod_{x,y:\Raw_N}\IC_N(x,y)\to \embed{x}=\embed{y},\\
  \mathsf{trunc} &:& \mathsf{isSet}(\Quot{\Raw_N}{\IC_N}).
\end{array}
\]
The semantic interface uses these principles in three roles.  First,
identity-criterion proofs become paths in the quotient.  Second, functions
between raw carriers that respect the relevant identity criteria factor as
maps between quotient common-noun types.  Third, set-valued observations and
projections descend when their raw definitions respect the identity criterion.

This interface is intentionally non-dependent.  It does not claim that every
dot-type projection, coercive subtyping rule, or dependent elimination
principle follows automatically.  Those stronger principles should be treated
as separate semantic or metatheoretic assumptions.

\section{Main Result and Proof Status}

For a common-noun setoid $C$ and a non-dependent set-valued target $B$, write
\[
  \mathsf{Obs}_{R}(C,B)
    = \sum_{f:\Raw_C\to B}\mathsf{respects}_C(f),
  \qquad
  \mathsf{Obs}_{Q}(C,B) = (\Quot{\Raw_C}{\IC_C}\to B).
\]
Precomposition with $\embed{-}$ defines
$\Phi:\mathsf{Obs}_{Q}(C,B)\to\mathsf{Obs}_{R}(C,B)$; quotient recursion
defines $\Psi$ in the other direction.

\begin{theorem}[Conditional converse representation]
Assume function extensionality, proof irrelevance for the respectfulness
field, and proposition-valued induction over the quotient presentation.
Then $\Phi\circ\Psi=\mathsf{id}$ and
$\Psi\circ\Phi=\mathsf{id}$.
\end{theorem}

\begin{proof}[Checked proof decomposition]
Every use of quotient recursion carries a setness witness for $B$.  The raw round trip follows pointwise from its beta rule;
function extensionality identifies the raw functions and proof irrelevance
identifies their respectfulness witnesses.  For the quotient round trip,
proposition-valued quotient induction reduces the equality to embedded
generators, where the same beta rule applies; function extensionality then
identifies the quotient observations.  Rocq checks the packaged result as
\path{conditional_observation_equivalence}.
\end{proof}

The main representation claim is conservative.  In the checked fragment, the
quotient-HIT presentation internalizes the identity criterion: identity proofs
in the setoid become paths in the quotient, and respectful set-valued
observations factor through quotient elimination.  The checked converse theorem is conditional:
the theorem above has exactly the three displayed premises.  No unconditional
converse axiom remains in the Rocq file.

\begin{table}[H]
\centering
\small
\caption{Submission claims and their current machine-checked status.}
\label{tab:proof-status}
\begin{tabularx}{\textwidth}{@{}>{\raggedright\arraybackslash}p{0.19\textwidth}>{\raggedright\arraybackslash}p{0.31\textwidth}X@{}}
\toprule
Claim & Rocq evidence & Status and boundary \\
\midrule
Identity internalization &
\path{passenger_identity_path}; \path{book_identity_needs_both_criteria} &
Checked from the quotient path constructor. \\
Respectful factorization &
\path{lifted_observation_adequacy}; book projection lemmas &
Checked from quotient recursion with an explicit set target, its beta rule, and respectfulness. \\
Converse representation &
\path{conditional_observation_equivalence} &
Checked under function extensionality, proof irrelevance, and quotient
proposition induction. \\
Unconditional converse & No Rocq declaration &
Not claimed; it remains a possible strengthening rather than an axiom. \\
Global assumptions & \path{AssumptionAudit.v} &
Exactly \path{qtype}, \path{embed}, \path{glue}, \path{qrec}, and
\path{qrec_beta}; all await native or library replacement. \\
\bottomrule
\end{tabularx}
\end{table}
The residual unconditional converse remains proof debt rather than a theorem
claimed by this draft.  This boundary is important for review.  The paper is
not a normalization or canonicity theorem for UTT with HITs; it is a semantic
adequacy claim for non-dependent set-valued observations out of quotient
common-noun types.

% Verification anchor: SCI-SUBMISSION-DRAFT-SKELETON.
% Verification anchor: SCI-DRAFT-MAIN-RESULT-SPINE.
% Verification anchor: SCI-CONVERSE-MAIN-CONDITIONAL.
% Verification anchors: SCI-CONVERSE-AXIOM-FREE-CONDITIONAL; SCI-QREC-SET-VALUED-TARGET; SCI-MAIN-THEOREM-ASSUMPTION-AUDIT.
% validation command: make -C work/formal proof-assistant-coq
\section{Linguistic Tests}

The following three tests make the semantic role of the quotient explicit.
Each begins with a raw carrier and an identity criterion; the quotient then
turns a proof of that criterion into a path.  Table~\ref{tab:case-artifacts}
records the corresponding checked lemmas.

\subsection{Passenger/person and person-times}

Let a passenger token record both a person and a journey:
\[
  \Raw_{\mathsf{Passenger}} =
  \mathsf{Person}\times\mathsf{Journey},
  \qquad
  (p,j)\;\IC_{\mathsf{Passenger}}\;(p',j')
  \quad\text{iff}\quad p=p'\ \wedge\ j=j'.
\]
For one person $p$ and two distinct journeys $j_1\neq j_2$, the raw tokens
$(p,j_1)$ and $(p,j_2)$ have the same person component but do not satisfy the
passenger identity criterion.  This is the intended person-times effect:
person equality alone is not a witness for passenger identity.  The Rocq lemma
\path{same_person_is_not_passenger_identity} checks this minimal countermodel,
while \path{passenger_identity_requires_journey} extracts journey equality
from any passenger-identity witness.  Conversely,
\path{passenger_identity_path} sends a full identity witness to equality in
the passenger quotient.  The example does not claim that the two quotient
terms are provably unequal; that would require an additional separation
principle.

Concretely, take one traveller with two distinct journey tokens, for example
two bus rides on different occasions.  The analysis preserves equality of the
traveller component while withholding a passenger-identity witness until the
journey components also agree.  It therefore models a person-times contrast
without inferring that the two tokens represent two persons.

\subsection{Modified common nouns}

For a modifier $M$ over a base noun $N$, the raw carrier records an individual
and modifier evidence:
\[
  \Raw_{M N}=\sum_{x:\Raw_N}M(x),
  \qquad
  (x,m)\;\IC_{M N}\;(y,n)
  \quad\text{iff}\quad x\;\IC_N\;y.
\]
Thus two proofs that the same individual is, for example, handsome do not
create two modified-noun instances merely because the proof terms differ.
\path{modifier_witnesses_do_not_count_twice} checks exactly that equality;
\path{base_identity_lifts_to_modified_identity} lifts base identity to the
modified noun.  The respectful projection back to the base noun is checked by
\path{modified_to_base_raw_respects} and computes on embedded tokens by
\path{modified_to_base_beta}.  This is a controlled treatment of modifier
evidence, not a general assertion that all semantic modifiers are
proof-irrelevant.

For instance, two independent derivations that the same individual is a
handsome man yield distinct raw evidence pairs but one modified-noun instance.
The case is about the individuation of the modified common noun; it does not
identify distinct individuals merely because they satisfy the same modifier.

\subsection{Book/copredication}

A raw book packages a physical and an informational component:
\[
  \Raw_{\mathsf{Book}}=\mathsf{Phys}\times\mathsf{Info},
  \qquad
  (p,i)\;\IC_{\mathsf{Book}}\;(p',i')
  \quad\text{iff}\quad
  p\;\IC_{\mathsf{Phys}}\;p'\ \wedge\ i\;\IC_{\mathsf{Info}}\;i'.
\]
Book identity is therefore conjunctive, while the two quotient observations
remain distinct: \path{book_to_physical_raw_respects} licenses the physical
projection and \path{book_to_informational_raw_respects} licenses the
informational one.  Their beta lemmas,
\path{book_to_physical_beta} and \path{book_to_informational_beta}, prove
that each quotient-level projection computes as the corresponding raw
projection on an embedded token.  This supports coordinated physical and
informational observations of one book without claiming that quotient HITs
alone provide a complete dot-type or coercion theory.

The familiar test sentence ``John picked up and mastered three books'' combines
physical predication with informational predication \cite{ChatzikyriakidisLuo2015}.
The quotient analysis does not derive the dot-type account itself.  It states
the additional identity condition under which the two projections concern the
same book tokens, and it checks that both projections descend through the
quotient interface.

\begin{table}[H]
\centering
\small
\caption{Linguistic tests and their checked Rocq anchors.}
\label{tab:case-artifacts}
\begin{tabularx}{\textwidth}{@{}>{\raggedright\arraybackslash}p{0.18\textwidth}>{\raggedright\arraybackslash}p{0.38\textwidth}X@{}}
\toprule
Case & Semantic distinction & Checked anchor \\
\midrule
Passenger & Same person is insufficient when journeys differ &
\scriptsize\texttt{same\_person\_is\_not\_passenger\_identity}\newline
\scriptsize\texttt{passenger\_identity\_path} \\
Modified CN & Distinct modifier witnesses do not multiply an instance &
\scriptsize\texttt{modifier\_witnesses\_do\_not\_count\_twice}\newline
\scriptsize\texttt{modified\_to\_base\_beta} \\
Book & Conjunctive book identity with two respectful projections &
\scriptsize\texttt{book\_to\_physical\_beta}\newline
\scriptsize\texttt{book\_to\_informational\_beta} \\
\bottomrule
\end{tabularx}
\end{table}

\section{Verification Status}

The current formal development is Coq/Rocq-first.  Plain Coq does not provide
native quotient HITs, so the checked file represents the semantic interface,
the respectfulness obligations, and the conditional converse theorem rather
than a native HIT implementation.  The current build checks the Coq-facing
entry point through the formal Makefile and records the remaining HIT-native
work as explicit proof debt.

For this submission, the required mechanisation gate is:
\begin{quote}
\small\texttt{make -C work/formal proof-assistant-sci}
\end{quote}
It first compiles the formal development with Rocq 9.0.1 and then checks the
exact global assumptions of the packaged main theorem.  Cubical Agda is not
part of that acceptance gate: a future Agda or HoTT port would be an
independent replication of the quotient interface, not a condition for calling
the present Rocq claims checked.

The Coq signature requires an \path{is_hset} witness for every \path{qrec}
target.  Rocq's \emph{Print Assumptions} command reports exactly five global
dependencies for the packaged theorem: \path{qtype}, \path{embed}, \path{glue},
\path{qrec}, and \path{qrec_beta}.  The three non-constructive principles remain
quantified theorem premises, and no inverse-law axiom is reported.  The five
dependencies remain interface parameters rather than a native implementation.

The submission claim should therefore be stated in two layers.  The checked
layer is the Rocq quotient-interface development and conditional converse for
non-dependent set-valued observations.  The future layer is a native quotient-
HIT realization, whether in a Rocq-compatible HoTT library or in another proof
assistant, in which the quotient type former, path constructor, set
truncation, and eliminator are supplied rather than axiomatized as an
interface.

\section{Limitations and Conclusion}

The present artifact is a focused paper fragment rather than a finished journal
submission.  The main remaining tasks are to expand the semantic comparison
with the CN-as-types and dot-type literature, explain the linguistic payoff of
each test in more detail, and keep engineering artifacts out of the main text.
A native HIT realization remains future work, not a condition for the claims
made here.

The central conclusion is nevertheless stable: quotient HITs give a precise
way to turn CN-as-setoids into CNs-as-quotient-HITs while preserving the
identity-criterion insight that motivates common nouns as types.

\section*{Reproducibility}

The anonymized source package contains the Rocq development and the manuscript
source.  The command \path{make -C work/formal proof-assistant-sci} compiles
the Coq-facing development and audits the global assumptions of the packaged
converse theorem; \path{make -C work/paper submission.pdf} builds this
manuscript.  The separate JoLLI title page and declarations remain author-
supplied material and are intentionally excluded from the anonymous review
copy.

Table~\ref{tab:artifact-map} maps each paper-facing claim to its executable
artifact.  All paths are relative to the root of the accompanying source
package.  A versioned public release is identified in the author materials and
editorial metadata.  Its author-identifying URL is intentionally omitted from
this anonymous copy to preserve blinded review.

\begin{table}[H]
\centering
\small
\caption{Executable evidence for the paper's formal claims.}
\label{tab:artifact-map}
\begin{tabularx}{\textwidth}{@{}>{\raggedright\arraybackslash}p{0.24\textwidth}>{\raggedright\arraybackslash}p{0.36\textwidth}X@{}}
\toprule
Paper claim & Artifact & Reproduction command \\
\midrule
Quotient interface; passenger, modified-CN, and book lemmas &
\texttt{CNQuotientHit.v} &
\texttt{make -C work/formal}\newline\texttt{proof-assistant-coq} \\
Conditional converse representation &
{\scriptsize\texttt{conditional\_observation\_}\newline\texttt{equivalence}} &
\texttt{make -C work/formal}\newline\texttt{proof-assistant-sci} \\
Exact global-assumption boundary &
\texttt{AssumptionAudit.v} &
\texttt{make -C work/formal}\newline\texttt{proof-assistant-coq-assumptions} \\
JoLLI manuscript requirements &
\texttt{submission.tex}; \texttt{preflight.sh} &
\texttt{make -C work/paper}\newline\texttt{jolli-preflight} \\
\bottomrule
\end{tabularx}
\end{table}

\appendix
\section{Selected Rocq Evidence}

This appendix is deliberately a readable extract, not a replacement for the
complete source release.  It fixes the exact formal shape of the common-noun
setoid, the postulated quotient-HIT interface, one negative passenger test,
and one computation rule.  The complete files and build commands are listed
in Table~\ref{tab:artifact-map}.

\subsection{Setoids and the quotient interface}

\begin{verbatim}
Record CNSetoid := {
  carrier : Type;
  ic : carrier -> carrier -> Prop;
  ic_refl : forall x, ic x x;
  ic_sym : forall x y, ic x y -> ic y x;
  ic_trans : forall x y z, ic x y -> ic y z -> ic x z
}.

Definition respects (C : CNSetoid) (B : Type)
  (f : carrier C -> B) : Prop :=
  forall x y, ic C x y -> f x = f y.

Module QuotientHIT.
  Parameter qtype : CNSetoid -> Type.
  Parameter embed : forall C, carrier C -> qtype C.
  Parameter glue : forall (C : CNSetoid) (x y : carrier C),
    ic C x y -> embed C x = embed C y.
  Parameter qrec : forall (C : CNSetoid) (B : Type),
    is_hset B -> forall f : carrier C -> B,
    respects C B f -> qtype C -> B.
  Parameter qrec_beta : forall (C : CNSetoid) (B : Type)
    (B_is_set : is_hset B) (f : carrier C -> B)
    (r : respects C B f) (x : carrier C),
    qrec C B B_is_set f r (embed C x) = f x.
End QuotientHIT.
\end{verbatim}

The interface is explicit about its status: in plain Rocq these are parameters,
not a claim that Rocq natively implements quotient HITs.  The assumption audit
below establishes that the packaged main theorem has no additional global
assumptions.

\subsection{Passenger counterexample and quotient computation}

\begin{verbatim}
Lemma same_person_is_not_passenger_identity :
  fst first_trip = fst second_trip /\
  ~ same_passenger first_trip second_trip.
Proof.
  split; [reflexivity | intros H].
  destruct H as [_ Hjourney]. discriminate Hjourney.
Qed.

Lemma modified_to_base_beta :
  forall u : RawModified,
    modified_to_base (embed ModifiedCN u) =
      modified_to_base_raw u.
Proof. intro u. apply qrec_beta. Qed.
\end{verbatim}

The first lemma formalizes the person-times distinction: a shared person
component does not supply passenger identity when journeys differ.  The second
is the point computation rule for the respectful projection from a modified
common noun to its base noun.  The book projection beta lemmas are analogous
and remain in the complete source.

\subsection{Audited boundary}

The following build target compiles the development and runs
\path{Print Assumptions} on
\path{conditional_observation_equivalence}.  Its required result is exactly
the five quotient-interface constants \path{qtype}, \path{embed},
\path{glue}, \path{qrec}, and \path{qrec_beta}.  Function extensionality,
proof irrelevance, and quotient proposition induction occur as explicit
premises of the conditional theorem, not as undeclared global axioms.
\begin{quote}
\small\texttt{make -C work/formal}\\
\texttt{proof-assistant-sci}
\end{quote}

\begin{thebibliography}{9}

\bibitem[Coquand et~al.(2018)Coquand, Huber, and M\"ortberg]{CoquandHuberMortberg2018}
T. Coquand, S. Huber, and A. M{\"o}rtberg.
\newblock On Higher Inductive Types in Cubical Type Theory.
\newblock In \emph{Proceedings of the 33rd Annual ACM/IEEE Symposium on Logic
in Computer Science (LICS 2018)}, 2018.
\newblock arXiv:1802.01170.

\bibitem[Univalent Foundations Program(2013)]{HoTTBook}
The Univalent Foundations Program.
\newblock \emph{Homotopy Type Theory: Univalent Foundations of Mathematics}.
\newblock Institute for Advanced Study, 2013.

\bibitem[Chatzikyriakidis and Luo(2015)]{ChatzikyriakidisLuo2015}
S. Chatzikyriakidis and Z. Luo.
\newblock Individuation Criteria, Dot-types and Copredication: A View from
Modern Type Theories.
\newblock In \emph{Proceedings of the 14th Meeting on the Mathematics of
Language (MoL 2015)}, Chicago, 2015.

\bibitem[Chatzikyriakidis and Luo(2018)]{ChatzikyriakidisLuo2018}
S. Chatzikyriakidis and Z. Luo.
\newblock Identity Criteria of Common Nouns and dot-types for Copredication.
\newblock \emph{Oslo Studies in Language}, 10(2):121--141, 2018.

\bibitem[Luo(2012)]{Luo2012CN}
Z. Luo.
\newblock Common Nouns as Types.
\newblock In D. B{\'e}chet and A. Dikovsky, editors, \emph{Logical Aspects of
Computational Linguistics (LACL 2012)}, LNCS 7351, pages 173--185. Springer,
2012.

\bibitem[Luo(2024)]{Luo2024HIT}
Z. Luo.
\newblock Common Nouns as Types: Higher Inductive Types in Type-Theoretical
Semantics.
\newblock Invited talk at Logic Colloquium 2024, Gothenburg, Sweden, June 2024.

\bibitem[Lumsdaine and Shulman(2020)]{LumsdaineShulman2017}
P. L. Lumsdaine and M. Shulman.
\newblock Semantics of Higher Inductive Types.
\newblock \emph{Math. Proc. Camb. Phil. Soc.}, 169:159--208, 2020.

\end{thebibliography}

\end{document}
